\section{Introduction}\label{introduction}

Does social media sentiment causally affect intraday stock price
movements, and does the answer depend on the type of firm? This paper
provides evidence for both questions using three Bloomberg panels: a
daily S\&P~500 panel over 2025--2026 with a robust set of confounders; an extended panel of the S\&P~500 from 2016--2026; and an extended Russell 3000 panel
over 2016--2026.  These are used to employ firm fixed-effects OLS, Double Machine
Learning (DoubleML), and Fama-MacBeth cross-sectional regressions to
estimate the effect of Twitter/X sentiment on stock returns across the
majority of the US market.

The question matters for three core reasons. Social media sentiment has
been widely studied, but as its proliferation has increased, its value
as a news source has degraded. The ability to measure its influence in
the market provides a meaningful indicator of its usefulness as an
information sharing tool. Secondly, as retail investor numbers have
grown, the number of investors with informal access to market data has
also grown. These traders are more likely to use platforms such as
Twitter, Reddit, and StockTwits for information than formal, expensive
platforms such as Bloomberg. This investor segment is also more active
in the small and mid-cap sections of the market. Lastly, the question of
Twitter's causal influence remains open. While much literature
documents that text-derived sentiment correlates with subsequent
returns, correlation is not causation. Firms that attract more social
media attention may differ systematically from those that do not;
sentiment itself may simply respond to price changes rather than cause
them. Resolving this requires identification strategies capable of
isolating the sentiment effect from confounders and from reverse
causality. Investigating these topics provides useful information on how
information diffuses through the market and why.

Other research has shown that the fraction of negative words in the Wall
Street Journal \emph{Abreast of the Market} column predicts next-day Dow
Jones returns and trading volume, with the effect reversing within a
week. In a study by Gu and Kurov (2020) Fama-MacBeth regressions have
been used to show that lagged Twitter sentiment predicts open-to-open
stock returns in a 2015--2017 sample, with the effect concentrated at a
one-day lag (a study replicated in this paper for both robustness and
trend analysis). These studies provide the empirical background for this
analysis. % [R4]: Informal register — "goes a bit further" and "I sought to make" are weak.
% Suggest: "This paper extends the analysis in two directions:" then enumerate.
This paper goes a bit further by examining whether effects are
heterogeneous across the market-capitalization distribution and by
substantially testing for reverse causality. Based on that gap, I sought
to make three contributions.

The first is methodological. By applying the PLR model of
Chernozhukov et al.\ (2018), combined with XGBoost as a nuisance
learner, I obtain an estimate of the sentiment effect that is
asymptotically unconfounded by high-dimensional covariates, regardless
of the functional form of their relationship with the treatment and
outcome. The DoubleML estimator achieves this by cross-fitting: nuisance
models for both \(E\lbrack\text{twitter\_sent} \mid X\rbrack\) and
\(E\lbrack\text{return} \mid X\rbrack\) are trained on held-out folds,
so the final coefficient is not contaminated by in-sample overfitting.
This is, to my knowledge, the first application of DoubleML to intraday
Twitter sentiment and stock returns.

The second contribution is the systematic documentation of a reverse
causality concern. I pair the primary DoubleML estimates with a
firm-level placebo test that regresses yesterday's return on today's
Twitter sentiment. Under the null of no reverse causation, this
coefficient should be indistinguishable from zero. It is not. The
estimated coefficient is \(+ 2.348\) (\(p < 0.01\)), substantially
larger in magnitude than any of the forward-looking estimates. This
result is inconsistent with clean causal transmission from sentiment to
prices and raises the possibility that Bloomberg's sentiment measure
largely reflects the market's own prior-day behavior rather than an
independent information signal.

The third contribution is the cap-group analysis of the Russell 3000. By
extending the Fama-MacBeth and Panel OLS specifications to approximately
1.97~million firm-day observations across the full Russell 3000 and
partitioning firms by time-averaged market capitalization, I document a
sharp split in sentiment effect: \(+ {0.375}^{***}\) for large-cap firms
(market cap \(>\)\$10B); \(+ {0.804}^{***}\) for mid-cap firms
(\$2B--\$10B); and \(+ {0.942}^{***}\) for small-cap firms (\(<\)\$2B).
An interaction specification also confirms that the effect is linearly
decreasing with firm size
(\(\text{twitter\_sent} \times \text{log\_mkt\_cap} = - {0.148}^{*}\))
across the entire 11 year panel, indicating that the smaller the firm,
the stronger Twitter sentiment's influence. This also informs the
discrepancy between the null result on the S\&P~500 and the significant
effects documented in other studies that include smaller firms,
indicating that Twitter sentiment is most informative precisely where
institutional information is sparsest.

The main findings can be summarized as follows. In the S\&P~500 firm
fixed-effects specifications (\(N \approx 160,000\) firm-day
observations), \texttt{twitter\_sent} is statistically insignificant
across all four treatment constructions. News sentiment enters with a
coefficient of \(- 1.056\) (\(p < 0.01\)) and is the dominant predictor
in the linear OLS framework. The DoubleML estimate for
\texttt{twitter\_sent} is \(- 0.924\) (\(p < 0.01\)), but the placebo
coefficient of \(+ 2.348\) makes causal attribution very weak. The
Fama-MacBeth coefficient on lagged Twitter sentiment is insignificant in
every calendar year from 2016 through 2025 for the S\&P~500 universe,
with no structural break or temporal trend. For large-cap S\&P~500
constituents, Twitter sentiment provides no statistically reliable
predictive signal in any year examined.

The Russell 3000 results stand in contrast. The full-sample Panel OLS
coefficient on \texttt{twitter\_sent} is \(+ 0.708\) (\(p < 0.01\))
across 1.97~million firm-day observations, positive and significant
where the S\&P~500 estimate was negative and insignificant. The
Fama-MacBeth series by cap group shows consistent significance for mid-
and small-cap segments across most years, with the strongest effects
concentrated in 2018--2020. A daily long-short strategy---long
top-decile, short bottom-decile sentiment firms within each cap
group---generates positive returns in most years, suggesting that even
where the average linear coefficient is modest, Twitter sentiment
produces exploitable cross-sectional dispersion in realized returns.
This last result is interesting however, because this strategy was also
shown to be effective for the S\&P 500.

The remainder of the paper proceeds as follows.
Section~\ref{literature-review} reviews the relevant
literature on sentiment and asset pricing, social media as an
information channel, and causal identification in high-frequency panels.
Section~\ref{data} describes the Bloomberg panels,
variable construction, the Russell 3000 data infrastructure, and sample
properties. Section~\ref{methodology} presents the
FE-OLS, DoubleML, Fama-MacBeth, and cap-group specifications with full
notation and identifying assumptions.
Section~\ref{results} reports point estimates,
robustness checks, the placebo and lag-structure tests, and the full
Russell 3000 cap-group analysis.
Section~\ref{discussion} interprets the central tension
between the DoubleML estimates and the reverse causality evidence, and
the cap-group gradient in the context of the information environment
hypothesis. Section~\ref{conclusion} concludes with
limitations and directions for future work, including a custom
entity-level sentiment tool that may isolate causally relevant signal
from the Bloomberg composite.
